\documentclass[11pt,a4paper]{article}

\usepackage[margin=1in]{geometry}
\usepackage{amsmath,amssymb,amsfonts,amsthm}
\usepackage{bm}
\usepackage{hyperref}
\usepackage{enumitem}
\usepackage{mathtools}

\hypersetup{
    colorlinks=true,
    linkcolor=blue,
    urlcolor=blue,
    pdftitle={Advanced Probability and Statistics},
    pdfauthor={Your Name}
}

\newtheorem{definition}{Definition}[section]
\newtheorem{theorem}{Theorem}[section]
\newtheorem{lemma}{Lemma}[section]
\newtheorem{proposition}{Proposition}[section]
\newtheorem{corollary}{Corollary}[section]
\newtheorem{example}{Example}[section]
\newtheorem{remark}{Remark}[section]

\numberwithin{equation}{section}

\title{Advanced Probability and Statistics\\[0.5em]
       \large Measure-Theoretic Probability, Martingales, Stochastic Processes,\\
       and Asymptotic Statistical Theory}
\author{Theodosios Dimitrasopoulos\footnote{Email: theodosios.novak@gmail.com}}
\date{\today}

\begin{document}

\maketitle
\tableofcontents
\newpage

%%%%%%%%%%%%%%%%%%%%%%%%%%%%%%%%%%%%%%%%%%%%%%%%%%%%%%%%%%%%
\section{How to Use This Course}

This document assumes familiarity with intermediate probability and statistics: random vectors, multivariate normal, basic convergence concepts, CLT, MLE, likelihood ratio tests, and linear models.%
\footnote{Content and level are comparable to first-year graduate courses in measure-theoretic probability, advanced probability, and mathematical statistics syllabi.[web:61][web:62][web:69][web:71][web:72]}

The aim is to develop a rigorous measure-theoretic foundation for probability and to connect it to modern statistical theory: martingales, Markov chains, Brownian motion and stochastic calculus, decision-theoretic estimation, and asymptotic likelihood theory. 
Each major definition or theorem is followed by commentary and at least one illustrative example, often with a modeling flavor (e.g., finance, time series, Monte Carlo).

A suggested path:
\begin{enumerate}
    \item Sections \ref{sec:measure-foundations}--\ref{sec:weak-conv}: measure, integration, convergence.
    \item Sections \ref{sec:independence-bc}--\ref{sec:martingales}: independence, Borel--Cantelli, martingales.
    \item Sections \ref{sec:markov}--\ref{sec:brownian}: Markov chains and Brownian motion.
    \item Sections \ref{sec:decision}--\ref{sec:likelihood-asymptotics}: advanced mathematical statistics and likelihood asymptotics.
    \item Sections \ref{sec:concentration}--\ref{sec:high-dim}: concentration, large deviations, and high-dimensional hints.
\end{enumerate}

\newpage

%%%%%%%%%%%%%%%%%%%%%%%%%%%%%%%%%%%%%%%%%%%%%%%%%%%%%%%%%%%%
\section{Measure-Theoretic Foundations of Probability}
\label{sec:measure-foundations}

\subsection{Measure Spaces and Probability Spaces}

\begin{definition}[Sigma-algebra]
Let $\Omega$ be a non-empty set. 
A collection $\mathcal{F} \subseteq 2^\Omega$ is a \textbf{sigma-algebra} if:
\begin{enumerate}[label=(\roman*)]
    \item $\Omega \in \mathcal{F}$.
    \item If $A \in \mathcal{F}$, then $A^c \in \mathcal{F}$.
    \item If $A_1,A_2,\dots \in \mathcal{F}$, then $\displaystyle \bigcup_{n=1}^\infty A_n \in \mathcal{F}$.
\end{enumerate}
\end{definition}

A sigma-algebra formalizes the idea of sets on which we can consistently define probabilities. 
Closure under complements and countable unions ensures we can talk about ``or'', ``and'', and limits of events in a mathematically stable way.

\begin{definition}[Measure and probability measure]
A function $\mu : \mathcal{F} \to [0,\infty]$ is a \textbf{measure} if:
\begin{enumerate}[label=(\roman*)]
    \item $\mu(\varnothing) = 0$.
    \item For disjoint $A_1,A_2,\dots \in \mathcal{F}$,
    \[
      \mu\Big(\bigcup_{n=1}^\infty A_n\Big)
      = \sum_{n=1}^\infty \mu(A_n).
    \]
\end{enumerate}
If in addition $\mu(\Omega)=1$, then $\mu$ is a \textbf{probability measure}.
\end{definition}

\begin{definition}[Probability space]
A \textbf{probability space} is a triple $(\Omega,\mathcal{F},P)$ where $\mathcal{F}$ is a sigma-algebra on $\Omega$ and $P$ is a probability measure on $(\Omega,\mathcal{F})$.
\end{definition}

\subsection{Random Variables as Measurable Maps}

\begin{definition}[Random variable]
A function $X : (\Omega,\mathcal{F}) \to (\mathbb{R},\mathcal{B})$ is a \textbf{random variable} if it is measurable, i.e.,
\[
\{ \omega \in \Omega : X(\omega) \le x \} \in \mathcal{F}
\quad\text{for all } x \in \mathbb{R},
\]
where $\mathcal{B}$ is the Borel sigma-algebra on $\mathbb{R}$.
\end{definition}

The measurability condition guarantees that events defined via $X$ (such as $\{X\le x\}$, $\{X\in A\}$) are valid events in $\mathcal{F}$. 
From this viewpoint, a random variable is just a measurable transformation of the underlying probability space.

\begin{definition}[Distribution of a random variable]
The \textbf{distribution} of a random variable $X$ is the probability measure $P_X$ on $(\mathbb{R},\mathcal{B})$ defined by
\[
P_X(B) = P\big(X^{-1}(B)\big)
= P(\{\omega : X(\omega) \in B\}),\quad B \in \mathcal{B}.
\]
\end{definition}

\begin{example}[Pushforward measure for a function of a random variable]
Let $X$ be a random variable on $(\Omega,\mathcal{F},P)$ and consider $Y = g(X)$ where $g:\mathbb{R}\to\mathbb{R}$ is Borel-measurable. 
Then $Y$ is a random variable, and its distribution $P_Y$ is the \emph{pushforward} of $P_X$ by $g$. 
This is the measure-theoretic formalization of ``applying a function to a random variable transforms its distribution''.

\end{example}

\newpage

%%%%%%%%%%%%%%%%%%%%%%%%%%%%%%%%%%%%%%%%%%%%%%%%%%%%%%%%%%%%
\section{Integration and Expectation}

\subsection{Lebesgue Integral}

Let $(\Omega,\mathcal{F},P)$ be a probability space and $X$ a non-negative random variable. 
The Lebesgue integral $\int X\,dP$ is defined via monotone approximations by simple functions.

\begin{definition}[Simple function]
A function $s:\Omega \to [0,\infty)$ is \textbf{simple} if it takes finitely many values:
\[
s(\omega) = \sum_{k=1}^n a_k \mathbf{1}_{A_k}(\omega),
\]
where $a_k \ge 0$ and $A_k \in \mathcal{F}$ are disjoint.
\end{definition}

For a simple function $s$, its integral is
\[
\int s\,dP = \sum_{k=1}^n a_k P(A_k).
\]

For a general non-negative random variable $X$, we approximate $X$ from below by an increasing sequence of simple functions $s_n \uparrow X$ and define
\[
\int X\,dP = \sup_{n} \int s_n\,dP.
\]

\begin{definition}[Expectation]
If $X$ is integrable (i.e., $\int |X|\,dP < \infty$), its \textbf{expectation} is
\[
E[X] = \int X\,dP.
\]
\end{definition}

This generalizes discrete sums and continuous Riemann integrals into one unified definition. 
Expectation is linear, monotone, and continuous under certain limits (monotone or dominated convergence).

\subsection{Monotone Convergence and Dominated Convergence}

\begin{theorem}[Monotone Convergence Theorem]
If $0 \le X_1 \le X_2 \le \dots$ and $X_n \uparrow X$ almost surely, then
\[
\lim_{n\to\infty} \int X_n\,dP = \int X\,dP.
\]
\end{theorem}

\begin{theorem}[Dominated Convergence Theorem]
Let $X_n$ be a sequence of random variables such that $X_n \to X$ almost surely and there exists an integrable $Y$ with $|X_n| \le Y$ almost surely for all $n$. 
Then $X$ is integrable and
\[
\lim_{n\to\infty} \int X_n\,dP = \int X\,dP.
\]
\end{theorem}

\begin{example}[Passing limit inside expectation for Monte Carlo estimators]
Let $X_n(\omega) = \tfrac{1}{n}\sum_{i=1}^n g(Z_i(\omega))$ where $Z_i$ are i.i.d. with $E[|g(Z_1)|]<\infty$. 
Then $X_n \to E[g(Z_1)]$ almost surely by the strong law. 
If $|g(Z_i)|$ are dominated by an integrable $Y$, the dominated convergence theorem implies
\[
\lim_{n\to\infty} E[X_n] = E\big[ E[g(Z_1)] \big] = E[g(Z_1)].
\]
So one can interchange limit and expectation; this legitimizes taking limits in Monte Carlo approximations.

\end{example}

\newpage

%%%%%%%%%%%%%%%%%%%%%%%%%%%%%%%%%%%%%%%%%%%%%%%%%%%%%%%%%%%%
\section{Modes of Convergence and Weak Convergence}
\label{sec:weak-conv}

\subsection{Almost Sure, In Probability, In $L^p$}

We recall and sharpen the different convergence notions for random variables $X_n$ to $X$.

\begin{itemize}
    \item \textbf{Almost sure convergence}: $X_n \to X$ almost surely if
    \[
      P(\{\omega : X_n(\omega)\to X(\omega)\}) = 1.
    \]
    \item \textbf{Convergence in probability}: $X_n \xrightarrow{P} X$ if for every $\varepsilon>0$,
    \[
      \lim_{n\to\infty} P(|X_n - X| > \varepsilon) = 0.
    \]
    \item \textbf{Convergence in $L^p$}: $X_n \to X$ in $L^p$ if
    \[
      \lim_{n\to\infty} E[|X_n - X|^p] = 0,\quad 1 \le p < \infty.
    \]
\end{itemize}

The implications are (for integrable variables):
\[
X_n \xrightarrow{L^p} X \Longrightarrow X_n \xrightarrow{P} X,
\]
and
\[
X_n \xrightarrow{\text{a.s.}} X \Longrightarrow X_n \xrightarrow{P} X.
\]
The converses generally fail, illustrating that these are genuinely different notions.

\subsection{Weak Convergence of Distributions}

\begin{definition}[Weak convergence]
A sequence of random variables $X_n$ converges in distribution (or weakly) to $X$, written $X_n \Rightarrow X$, if
\[
\lim_{n\to\infty} F_{X_n}(x) = F_X(x)
\]
for all $x$ at which $F_X$ is continuous.
\end{definition}

Weak convergence is about distributions, not sample paths. 
It is the right notion for CLT-type results, large deviations, and asymptotic distributions of estimators.

\begin{theorem}[Continuous Mapping Theorem]
If $X_n \Rightarrow X$ and $g$ is continuous, then $g(X_n) \Rightarrow g(X)$.
\end{theorem}

\begin{example}[Delta method revisited]
If $\sqrt{n}(T_n-\theta) \Rightarrow N(0,\tau^2)$ and $g$ is differentiable at $\theta$ with $g'(\theta)\ne 0$, then by a Taylor expansion and the continuous mapping theorem,
\[
\sqrt{n}\big(g(T_n)-g(\theta)\big) \Rightarrow N\big(0,(g'(\theta))^2\tau^2\big).
\]
This formalizes the intuitive idea that small perturbations in $T_n$ are linearly amplified by $g$ in the limit.

\end{example}

\newpage

%%%%%%%%%%%%%%%%%%%%%%%%%%%%%%%%%%%%%%%%%%%%%%%%%%%%%%%%%%%%
\section{Independence, Borel--Cantelli, and Zero--One Laws}
\label{sec:independence-bc}

\subsection{Independence of Sigma-Algebras and Random Variables}

\begin{definition}[Independence of sigma-algebras]
Sigma-algebras $\mathcal{F}_1,\dots,\mathcal{F}_n$ are \textbf{independent} if
\[
P(A_1\cap\cdots\cap A_n) = \prod_{i=1}^n P(A_i)
\]
for all $A_i\in \mathcal{F}_i$.
\end{definition}

Random variables $X_i$ are independent if the sigma-algebras they generate are independent. 
This abstraction is crucial for product probability spaces and infinite collections of variables.

\subsection{Borel--Cantelli Lemmas}

\begin{theorem}[First Borel--Cantelli Lemma]
If $A_1,A_2,\dots$ is a sequence of events with
\[
\sum_{n=1}^\infty P(A_n) < \infty,
\]
then
\[
P(A_n \text{ i.o.}) = 0,
\]
where ``i.o.'' means ``infinitely often''.
\end{theorem}

This says that if the total expected number of occurrences is finite, then with probability 1, only finitely many occur.

\begin{theorem}[Second Borel--Cantelli Lemma]
If $A_1,A_2,\dots$ are independent events and
\[
\sum_{n=1}^\infty P(A_n) = \infty,
\]
then
\[
P(A_n \text{ i.o.}) = 1.
\]
\end{theorem}

\begin{example}[Gambler eventually goes broke?]
Let $A_n$ be the event that a gambler's fortune hits zero on day $n$. 
If the $A_n$ were independent and the sum of their probabilities diverged, Borel--Cantelli would say the gambler will go broke infinitely often with probability 1. 
In realistic models, the $A_n$ are not independent, but Borel--Cantelli still provides insight: if ruin probabilities per day decay slowly enough, ruin becomes almost certain in the long run.

\end{example}

\subsection{Kolmogorov Zero--One Law}

\begin{definition}[Tail sigma-algebra]
For a sequence of random variables $(X_n)$, the \textbf{tail sigma-algebra} is
\[
\mathcal{T} = \bigcap_{n=1}^\infty \sigma(X_k : k\ge n).
\]
\end{definition}

\begin{theorem}[Kolmogorov zero--one law]
If $(X_n)$ are independent random variables, then any event $A \in \mathcal{T}$ has probability $0$ or $1$.
\end{theorem}

Tail events depend only on what happens ``in the long run''. 
The zero--one law says there is no meaningful randomness left at infinity: such events either almost surely happen or almost surely do not.

\begin{example}[Convergence events]
For independent $(X_n)$, the event $\{\limsup X_n = \infty\}$ is a tail event. 
Hence it has probability 0 or 1. 
Determining which requires additional information, but the zero--one law tells us there are no intermediate probabilities for such long-run behaviors.

\end{example}

\newpage

%%%%%%%%%%%%%%%%%%%%%%%%%%%%%%%%%%%%%%%%%%%%%%%%%%%%%%%%%%%%
\section{Strong Laws, Central Limit Theorems, and LIL}

\subsection{Kolmogorov's Strong Law of Large Numbers}

Let $(X_n)$ be independent with $E[X_n]=0$ and $\sum_{n=1}^\infty \operatorname{Var}(X_n)/n^2 < \infty$. 
Then
\[
\frac{1}{n}\sum_{k=1}^n X_k \xrightarrow{\text{a.s.}} 0.
\]

For i.i.d.\ $X_n$ with $E[|X_1|]<\infty$, we get the classical strong law:
\[
\bar{X}_n \xrightarrow{\text{a.s.}} E[X_1].
\]

\subsection{Lindeberg--Feller CLT}

\begin{theorem}[Lindeberg--Feller CLT, simplified]
Let $X_{n,k}$ be a triangular array of independent mean-zero random variables with variances $\sigma_{n,k}^2$, and define
\[
s_n^2 = \sum_{k=1}^{k_n} \sigma_{n,k}^2.
\]
If $s_n^2 \to \sigma^2 > 0$ and Lindeberg's condition holds:
\[
\forall \varepsilon>0,\quad 
\lim_{n\to\infty} \frac{1}{s_n^2}\sum_{k=1}^{k_n} E\big[
X_{n,k}^2 \mathbf{1}_{\{|X_{n,k}|>\varepsilon s_n\}}
\big] = 0,
\]
then
\[
\frac{1}{s_n}\sum_{k=1}^{k_n} X_{n,k} \Rightarrow N(0,\sigma^2).
\]
\end{theorem}

This generalizes the classical CLT to independent but not identically distributed variables, as long as no single term dominates the variance and the Lindeberg condition controls large jumps.

\subsection{Law of the Iterated Logarithm (LIL)}

For i.i.d.\ mean-zero variance-one variables $X_n$, define partial sums $S_n = \sum_{k=1}^n X_k$. 
The LIL states that almost surely,
\[
\limsup_{n\to\infty} \frac{S_n}{\sqrt{2n\log\log n}} = 1,\quad
\liminf_{n\to\infty} \frac{S_n}{\sqrt{2n\log\log n}} = -1.
\]

This describes the ultimate envelope of fluctuations of $S_n$: although $S_n$ is of order $\sqrt{n}$, the $\sqrt{2n\log\log n}$ normalization captures the extreme oscillations infinitely often.

\begin{example}[Random walk extremes]
Consider an unbiased random walk $S_n$ on $\mathbb{Z}$ with step $\pm 1$. 
The LIL tells you that $S_n$ visits arbitrarily large values of the order $\sqrt{2n\log\log n}$ infinitely often, but almost never significantly larger relative to this scale. 
For finance, it hints at how extreme deviations from expected P\&L can scale with horizon when modeling returns as sums of independent shocks.

\end{example}

\newpage

%%%%%%%%%%%%%%%%%%%%%%%%%%%%%%%%%%%%%%%%%%%%%%%%%%%%%%%%%%%%
\section{Conditional Expectation as Projection and Martingales}
\label{sec:martingales}

\subsection{Conditional Expectation w.r.t.\ a Sigma-Algebra}

\begin{definition}[Conditional expectation]
Let $X$ be integrable and $\mathcal{G} \subseteq \mathcal{F}$ a sub-sigma-algebra. 
A random variable $Y$ is the \textbf{conditional expectation} of $X$ given $\mathcal{G}$, written $Y = E[X \mid \mathcal{G}]$, if:
\begin{enumerate}[label=(\roman*)]
    \item $Y$ is $\mathcal{G}$-measurable.
    \item For all $G \in \mathcal{G}$,
    \[
    \int_G Y\,dP = \int_G X\,dP.
    \]
\end{enumerate}
\end{definition}

Intuitively, $E[X \mid \mathcal{G}]$ is the best $\mathcal{G}$-measurable approximation to $X$ in $L^2$ sense; it is the orthogonal projection of $X$ onto the closed subspace of $\mathcal{G}$-measurable square-integrable random variables.

\subsection{Martingales}

\begin{definition}[Martingale]
Let $(\mathcal{F}_n)$ be a filtration, i.e., an increasing sequence of sigma-algebras. 
A sequence $(M_n)$ of integrable random variables is a \textbf{martingale} w.r.t.\ $(\mathcal{F}_n)$ if:
\begin{enumerate}[label=(\roman*)]
    \item $M_n$ is $\mathcal{F}_n$-measurable.
    \item $E[|M_n|]<\infty$.
    \item $E[M_{n+1} \mid \mathcal{F}_n] = M_n$ almost surely.
\end{enumerate}
\end{definition}

A martingale models a ``fair game'': the conditional expectation of tomorrow's value given today's information is today's value.

\begin{example}[Discounted asset price under risk-neutral measure]
In arbitrage-free pricing, under a risk-neutral measure $Q$ and constant interest rate $r$, the discounted asset price process
\[
M_n = e^{-rn} S_n
\]
is typically a martingale w.r.t.\ the natural filtration of $(S_n)$. 
This encodes the idea that, under risk-neutral probabilities, the expected discounted price tomorrow equals today's discounted price.

\end{example}

\subsection{Optional Stopping and Doob's Inequalities}

\begin{definition}[Stopping time]
A random time $\tau:\Omega\to\{0,1,2,\dots,\infty\}$ is a \textbf{stopping time} w.r.t.\ $(\mathcal{F}_n)$ if $\{\tau \le n\}\in \mathcal{F}_n$ for all $n$.
\end{definition}

\begin{theorem}[Optional Stopping (simplified)]
Let $(M_n)$ be a martingale and $\tau$ a bounded stopping time (i.e., $\tau \le N$ almost surely). 
Then
\[
E[M_\tau] = E[M_0].
\]
\end{theorem}

This formalizes the idea that you cannot beat a fair game by deciding when to stop, as long as your stopping rule is based on current information and bounded.

\begin{theorem}[Doob's $L^p$ Inequality (simplified)]
If $(M_n)$ is a martingale with $M_0=0$ and $1<p<\infty$, then there exists $C_p$ such that
\[
E\big[\max_{1\le k\le n} |M_k|^p\big]
\le C_p\, E[|M_n|^p].
\]
\end{theorem}

This inequality controls the maximal fluctuations of a martingale in terms of its terminal value, a crucial tool in proving convergence and limit theorems.

\subsection{Martingale Convergence}

\begin{theorem}[Martingale Convergence Theorem, one version]
Let $(M_n)$ be a martingale that is bounded in $L^1$ (i.e., $\sup_n E[|M_n|]<\infty$). 
Then there exists an integrable random variable $M_\infty$ such that
\[
M_n \xrightarrow{\text{a.s.}} M_\infty
\quad\text{and}\quad
M_n \xrightarrow{L^1} M_\infty.
\]
\end{theorem}

This says fair games with uniformly bounded expected absolute value converge almost surely to a limiting random variable; in finance, it relates to terminal values under risk-neutral measures.

\newpage

%%%%%%%%%%%%%%%%%%%%%%%%%%%%%%%%%%%%%%%%%%%%%%%%%%%%%%%%%%%%
\section{Discrete-Time Markov Chains}
\label{sec:markov}

\subsection{Definitions}

\begin{definition}[Markov chain]
A sequence $(X_n)_{n\ge 0}$ of random variables on a countable state space $S$ is a \textbf{Markov chain} with transition probabilities $P=(p_{ij})$ if
\[
P(X_{n+1}=j \mid X_n=i, X_{n-1}=i_{n-1},\dots,X_0=i_0)
= P(X_{n+1}=j \mid X_n=i) = p_{ij},
\]
for all states $i_0,\dots,i_{n-1},i,j\in S$.
\end{definition}

This expresses the \emph{Markov property}: the future depends on the past only through the present.

\subsection{Transition Probabilities and $n$-Step Transitions}

The $n$-step transition probabilities are
\[
p_{ij}^{(n)} = P(X_n=j \mid X_0=i).
\]

These can be computed via matrix powers:
\[
P^{(n)} = P^n,\quad\text{where } (P^n)_{ij} = p_{ij}^{(n)}.
\]

\subsection{Classification of States and Stationary Distributions}

States can be recurrent or transient, periodic or aperiodic, and the chain may be irreducible (every state communicates with every other). 
A probability vector $\pi$ on $S$ is \textbf{stationary} if
\[
\pi^\top P = \pi^\top.
\]

\begin{example}[Random walk on integers mod $m$]
Let $S=\{0,1,\dots,m-1\}$ and
\[
P(X_{n+1} = i+1 \mod m \mid X_n=i) = 1.
\]
Then the chain just cycles deterministically through the states; every state is visited, but the chain is periodic with period $m$. 
A stationary distribution exists (uniform on $S$), but the chain does not converge to it from all starting states because of periodicity.

\end{example}

\subsection{Ergodic Theorem for Markov Chains}

Under irreducibility, aperiodicity, and positive recurrence, a unique stationary distribution $\pi$ exists and
\[
\lim_{n\to\infty} p_{ij}^{(n)} = \pi_j,\quad \forall i,j.
\]
Moreover, time-averages converge almost surely:
\[
\frac{1}{n}\sum_{k=1}^n f(X_k) \xrightarrow{\text{a.s.}} \sum_{j\in S} f(j)\,\pi_j,
\]
for any bounded function $f$. 
This is an ergodic theorem: sample paths reveal stationary expectations in the long run.

\begin{example}[MCMC intuition]
In Markov chain Monte Carlo, you construct a Markov chain whose stationary distribution is the target $\pi$. 
The ergodic theorem justifies using long-run averages of $f(X_k)$ along the chain to estimate $E_\pi[f(X)]$; the mathematical conditions ensure convergence of these averages.

\end{example}

\newpage

%%%%%%%%%%%%%%%%%%%%%%%%%%%%%%%%%%%%%%%%%%%%%%%%%%%%%%%%%%%%
\section{Brownian Motion and Introduction to Stochastic Calculus}
\label{sec:brownian}

\subsection{Standard Brownian Motion}

\begin{definition}[Standard Brownian motion]
A process $(B_t)_{t\ge 0}$ is a \textbf{standard Brownian motion} if:
\begin{enumerate}[label=(\roman*)]
    \item $B_0=0$ almost surely.
    \item $B_t$ has independent increments.
    \item For $0\le s < t$, $B_t - B_s \sim N(0,t-s)$.
    \item $B_t$ has continuous sample paths almost surely.
\end{enumerate}
\end{definition}

Brownian motion is the canonical continuous-time Gaussian process with stationary independent increments. 
It models continuous-time random noise and underlies diffusion processes, option pricing, and many physical phenomena.

\subsection{Quadratic Variation}

For a partition $\Pi_n = \{0=t_0 < t_1 < \cdots < t_n = t\}$, define
\[
V_t(\Pi_n) = \sum_{k=0}^{n-1} (B_{t_{k+1}} - B_{t_k})^2.
\]
As the mesh of the partition tends to 0, $V_t(\Pi_n)$ converges in probability to $t$. 
We call $[B]_t = t$ the \textbf{quadratic variation} of Brownian motion.

Quadratic variation captures the ``roughness'' of Brownian paths: unlike smooth functions, whose quadratic variation is zero, Brownian paths accumulate non-trivial variance along every interval.

\subsection{It\^o Integral (Idea)}

For an adapted process $(H_t)$ satisfying suitable integrability, the It\^o integral
\[
\int_0^t H_s\,dB_s
\]
is defined as an $L^2$-limit of stochastic Riemann sums
\[
\sum_{k} H_{t_k} (B_{t_{k+1}} - B_{t_k}).
\]

Key properties include:
\begin{itemize}
    \item Linearity in $H$.
    \item Isometry: $E\left[ \left( \int_0^t H_s\,dB_s \right)^2 \right] = E\left[ \int_0^t H_s^2\,ds \right]$.
\end{itemize}

\subsection{It\^o's Formula}

\begin{theorem}[It\^o's Formula, one dimension]
Let $X_t$ solve
\[
dX_t = \mu(t,X_t)\,dt + \sigma(t,X_t)\,dB_t,
\]
and let $f(t,x)$ be $C^{1,2}$ (once differentiable in $t$, twice in $x$). 
Then
\[
df(t,X_t)
= \left( f_t(t,X_t) + \mu(t,X_t) f_x(t,X_t)
+ \frac{1}{2}\sigma^2(t,X_t) f_{xx}(t,X_t) \right) dt
+ \sigma(t,X_t) f_x(t,X_t)\,dB_t.
\]
\end{theorem}

This is the stochastic chain rule: extra term $\tfrac{1}{2}\sigma^2 f_{xx}$ arises from quadratic variation. 
In quantitative finance, applying It\^o's formula to option prices leads to PDEs like Black--Scholes.

\begin{example}[Geometric Brownian motion]
Consider the SDE
\[
dS_t = \mu S_t\,dt + \sigma S_t\,dB_t,
\]
with $S_0>0$. 
Let $f(x) = \log x$, so $f'(x) = 1/x$, $f''(x) = -1/x^2$. 
Applying It\^o to $f(S_t)$ gives
\[
d\log S_t
= \left(\mu - \frac{1}{2}\sigma^2\right) dt + \sigma\,dB_t.
\]
Integrating yields
\[
S_t = S_0\exp\left(\Big(\mu-\frac{1}{2}\sigma^2\Big)t + \sigma B_t\right).
\]
Thus geometric Brownian motion has log-normal marginals; this is the classical continuous-time model for asset prices in Black--Scholes.

\end{example}

\newpage

%%%%%%%%%%%%%%%%%%%%%%%%%%%%%%%%%%%%%%%%%%%%%%%%%%%%%%%%%%%%
\section{Decision Theory and Risk in Statistical Inference}
\label{sec:decision}

\subsection{Decision Functions and Risk}

In mathematical statistics, an estimator is viewed as a decision rule under uncertainty.

\begin{definition}[Decision rule and loss]
Let $\Theta$ be the parameter space and $\mathcal{A}$ an action space. 
A \textbf{decision rule} $\delta$ maps the sample space to actions.
A \textbf{loss function} $L(\theta,a)$ quantifies the cost of taking action $a$ when the true parameter is $\theta$.
\end{definition}

\begin{definition}[Risk function]
The \textbf{risk} of a decision rule $\delta$ at $\theta$ is
\[
R(\theta,\delta) = E_\theta[L(\theta,\delta(X))],
\]
the expected loss under data generated from $\theta$.
\end{definition}

\subsection{Admissibility and Minimaxity}

A rule $\delta_1$ \emph{dominates} $\delta_2$ if $R(\theta,\delta_1) \le R(\theta,\delta_2)$ for all $\theta$, with strict inequality for some $\theta$. 
A rule is \textbf{admissible} if no other rule dominates it.

\begin{definition}[Minimax rule]
A rule $\delta^*$ is \textbf{minimax} if it minimizes the maximum risk:
\[
\sup_{\theta} R(\theta,\delta^*)
= \inf_{\delta} \sup_\theta R(\theta,\delta).
\]
\end{definition}

Minimax rules are conservative: they are optimal under worst-case parameter choices.

\subsection{Bayes Rules}

Given a prior $\pi$ on $\Theta$, the \textbf{Bayes risk} of $\delta$ is
\[
r(\pi,\delta)
= \int R(\theta,\delta)\,\pi(d\theta).
\]

A \textbf{Bayes rule} $\delta_\pi$ minimizes $r(\pi,\delta)$ over all $\delta$. 
Under regularity, Bayes rules are often admissible and have appealing decision-theoretic properties.

\begin{example}[Quadratic loss and posterior mean]
For estimating a scalar $\theta$ with squared error loss $L(\theta,a)=(\theta-a)^2$, the Bayes rule is the posterior mean $E[\theta\mid X]$. 
Indeed, for any posterior distribution, the value $a$ minimizing $E[(\theta-a)^2 \mid X]$ is $a=E[\theta\mid X]$, an elementary calculus fact.

\end{example}

\newpage

%%%%%%%%%%%%%%%%%%%%%%%%%%%%%%%%%%%%%%%%%%%%%%%%%%%%%%%%%%%%
\section{Sufficiency, Completeness, and UMVU Estimation}

\subsection{Sufficiency and Minimal Sufficiency}

We recall sufficiency:

\begin{definition}[Sufficient statistic]
A statistic $T(X)$ is \textbf{sufficient} for $\theta$ if the conditional distribution of $X$ given $T(X)$ does not depend on $\theta$.
\end{definition}

\begin{definition}[Minimal sufficient]
A sufficient statistic $T$ is \textbf{minimal sufficient} if it is a function of any other sufficient statistic: if $S$ is sufficient, then there exists $h$ such that $T=h(S)$ almost surely.
\end{definition}

In exponential families, certain natural statistics (like sums, sums of squares) are minimal sufficient.

\subsection{Completeness and Lehmann--Scheff\'e}

\begin{definition}[Complete statistic]
A statistic $T$ is \textbf{complete} for a family $\{P_\theta\}$ if for any measurable function $g$,
\[
E_\theta[g(T)] = 0 \text{ for all } \theta
\quad\Rightarrow\quad
P_\theta(g(T)=0)=1 \text{ for all } \theta.
\]
\end{definition}

\begin{theorem}[Lehmann--Scheff\'e]
If $T$ is a complete sufficient statistic for $\theta$, and $\hat{\phi}(T)$ is an unbiased estimator of $\phi(\theta)$, then $\hat{\phi}(T)$ is the unique uniformly minimum variance unbiased (UMVU) estimator of $\phi(\theta)$.
\end{theorem}

\begin{example}[Normal mean with known variance]
Let $X_1,\dots,X_n \sim N(\mu,\sigma^2)$ with known $\sigma^2$. 
The sample mean $\bar{X}$ is sufficient and complete for $\mu$ in this exponential family. 
Any unbiased estimator of $\mu$ that is a function of $\bar{X}$ must be $\bar{X}$ itself. 
Hence $\bar{X}$ is UMVU for $\mu$.

\end{example}

\subsection{Rao--Blackwell Theorem}

\begin{theorem}[Rao--Blackwell]
Let $T$ be sufficient for $\theta$ and let $\delta(X)$ be any estimator of $\phi(\theta)$ with finite risk under squared loss. 
Define the improved estimator
\[
\delta^*(X) = E_\theta[\delta(X) \mid T].
\]
Then $R(\theta,\delta^*) \le R(\theta,\delta)$ for all $\theta$. 
If $T$ is complete and $\delta$ is unbiased, then $\delta^*$ is UMVU.
\end{theorem}

Intuitively, conditioning on $T$ discards irrelevant noise while retaining all information about $\theta$, reducing variance without introducing bias.

\newpage

%%%%%%%%%%%%%%%%%%%%%%%%%%%%%%%%%%%%%%%%%%%%%%%%%%%%%%%%%%%%
\section{Information Inequalities and Cram\'er--Rao Bounds}

\subsection{Fisher Information}

For a parametric family $\{f(x;\theta)\}$ (dominated by a common measure), the \textbf{score} is
\[
U_\theta(X) = \frac{\partial}{\partial\theta}\log f(X;\theta),
\]
assuming differentiability under the integral sign.

\begin{definition}[Fisher information]
The \textbf{Fisher information} in one observation is
\[
I(\theta) = E_\theta[U_\theta(X)^2]
= -E_\theta\left[\frac{\partial^2}{\partial\theta^2}\log f(X;\theta)\right],
\]
when the interchange of derivative and expectation is valid.
\end{definition}

For i.i.d.\ $X_1,\dots,X_n$, the information in the sample is $n I(\theta)$.

\subsection{Cram\'er--Rao Inequality}

\begin{theorem}[Cram\'er--Rao lower bound]
Let $T(X)$ be an unbiased estimator of $\phi(\theta)$ with finite variance. 
Under regularity conditions,
\[
\operatorname{Var}_\theta(T)
\ge \frac{(\phi'(\theta))^2}{I(\theta)}.
\]
For i.i.d.\ samples, replace $I(\theta)$ by $nI(\theta)$.
\end{theorem}

This provides a lower bound on the variance of any unbiased estimator; estimators attaining the bound (for all $\theta$) are called \emph{efficient}.

\begin{example}[Normal mean]
Let $X_1,\dots,X_n \sim N(\mu,\sigma^2)$ with known $\sigma^2$, and consider estimating $\phi(\mu)=\mu$. 
The log-likelihood for one observation is
\[
\log f(x;\mu) = -\frac{1}{2}\log(2\pi\sigma^2) - \frac{(x-\mu)^2}{2\sigma^2},
\]
so
\[
U_\mu(X) = \frac{\partial}{\partial\mu}\log f(X;\mu)
= \frac{X-\mu}{\sigma^2},
\]
implying $I(\mu) = 1/\sigma^2$. 
For $n$ observations, $I_n(\mu) = n/\sigma^2$, so the CRLB for any unbiased estimator $T$ of $\mu$ is
\[
\operatorname{Var}(T) \ge \frac{1}{I_n(\mu)} = \frac{\sigma^2}{n}.
\]
The sample mean $\bar{X}$ achieves this bound, so it is efficient.

\end{example}

\newpage

%%%%%%%%%%%%%%%%%%%%%%%%%%%%%%%%%%%%%%%%%%%%%%%%%%%%%%%%%%%%
\section{Asymptotic Likelihood Theory}
\label{sec:likelihood-asymptotics}

\subsection{Consistency and Asymptotic Normality of MLEs}

Under standard regularity conditions (identifiability, differentiability, integrability, etc.), the maximum likelihood estimator $\hat{\theta}_n$ of a scalar parameter $\theta_0$ is consistent and asymptotically normal:
\[
\hat{\theta}_n \xrightarrow{P} \theta_0,
\]
\[
\sqrt{n}(\hat{\theta}_n - \theta_0)
\Rightarrow N\left(0,\,\frac{1}{I(\theta_0)}\right),
\]
where $I(\theta_0)$ is the Fisher information for one observation.

In vector form, if $\theta\in\mathbb{R}^p$, then
\[
\sqrt{n}(\hat{\theta}_n - \theta_0)
\Rightarrow N_p\left(\bm{0},\, I(\theta_0)^{-1}\right),
\]
with $I(\theta_0)$ now the Fisher information matrix.

\subsection{Likelihood Ratio, Score, and Wald Tests}

For testing $H_0:\theta=\theta_0$ vs.\ $H_1:\theta\ne\theta_0$, three asymptotically equivalent test statistics are:

\begin{itemize}
    \item \textbf{Likelihood ratio test} (LRT):
    \[
    \Lambda = \frac{L(\theta_0)}{L(\hat{\theta}_n)},\quad
    -2\log\Lambda
    \overset{H_0}{\Rightarrow} \chi^2_1.
    \]
    \item \textbf{Score test}:
    \[
    S = \frac{U_n(\theta_0)^2}{I_n(\theta_0)}
    \overset{H_0}{\Rightarrow} \chi^2_1,
    \]
    where $U_n(\theta)=\partial \ell_n(\theta)/\partial\theta$ is the sample score.
    \item \textbf{Wald test}:
    \[
    W = n\frac{(\hat{\theta}_n - \theta_0)^2}{\hat{V}_n}
    \overset{H_0}{\Rightarrow} \chi^2_1,
    \]
    where $\hat{V}_n$ is a consistent estimator of $\operatorname{Var}(\hat{\theta}_n)$.
\end{itemize}

These statistics all asymptotically follow the same chi-square distribution under $H_0$, though their finite-sample behavior and implementation details differ.

\begin{example}[Logistic regression coefficient test]
In logistic regression, suppose we test whether a particular coefficient $\beta_j$ equals zero. 
The Wald test uses $\hat{\beta}_j$ and its estimated standard error; the score test uses the derivative of the log-likelihood at $\beta_j=0$; the LRT compares maximized log-likelihoods with and without the predictor. 
Asymptotically, all three tests have size $\alpha$ and similar power properties, but in small samples they can behave differently.

\end{example}

\newpage

%%%%%%%%%%%%%%%%%%%%%%%%%%%%%%%%%%%%%%%%%%%%%%%%%%%%%%%%%%%%
\section{Bayesian Asymptotics and Bernstein--von Mises}

\subsection{Posterior Concentration}

Under regularity, Bayesian posteriors tend to concentrate around the true parameter as $n\to\infty$. 
In many models, the posterior variance shrinks at rate $1/n$ and the posterior mean is consistent for the true parameter.

\subsection{Bernstein--von Mises Theorem (Informal)}

Roughly, in regular parametric models with prior density $\pi(\theta)$ that is positive and continuous near the true parameter $\theta_0$, the posterior distribution of $\theta$ is asymptotically normal:
\[
\pi(\theta \mid X_1,\dots,X_n)
\approx N\left(\hat{\theta}_n,\, [n I(\theta_0)]^{-1}\right),
\]
where $\hat{\theta}_n$ is the MLE. 
More precisely, the posterior and the normal approximation are close in total variation.

This means that, asymptotically, Bayesian credible intervals and frequentist confidence intervals coincide in regular problems.

\begin{example}[Normal mean with conjugate prior]
Let $X_i\sim N(\mu,\sigma^2)$ with known $\sigma^2$, prior $\mu \sim N(\mu_0,\tau^2)$. 
The posterior is
\[
\mu \mid X_1,\dots,X_n \sim
N\left(
\frac{\mu_0/\tau^2 + n\bar{X}/\sigma^2}{1/\tau^2 + n/\sigma^2},
\ \frac{1}{1/\tau^2 + n/\sigma^2}
\right).
\]
As $n\to\infty$, the posterior mean approaches $\bar{X}$ and the posterior variance approaches $\sigma^2/n$, matching the frequentist asymptotic distribution of $\bar{X}$. 
This is a concrete illustration of Bernstein--von Mises in a simple setting.

\end{example}

\newpage

%%%%%%%%%%%%%%%%%%%%%%%%%%%%%%%%%%%%%%%%%%%%%%%%%%%%%%%%%%%%
\section{Concentration Inequalities and Large Deviations}
\label{sec:concentration}

\subsection{Hoeffding's Inequality}

\begin{theorem}[Hoeffding]
Let $X_1,\dots,X_n$ be independent with $a_i \le X_i \le b_i$ almost surely and $E[X_i]=\mu_i$. 
Then for $S_n = \sum_{i=1}^n X_i$ and any $t>0$,
\[
P(S_n - E[S_n] \ge t)
\le \exp\left(-\frac{2t^2}{\sum_{i=1}^n (b_i-a_i)^2}\right).
\]
\end{theorem}

This inequality shows that the sum of bounded independent variables is sharply concentrated around its mean, with exponential tails. 
Such results are fundamental in learning theory, empirical process theory, and high-dimensional statistics.[web:65]

\subsection{Chernoff Bounds}

Chernoff bounds use moment generating functions to produce exponential tail bounds. 
If $X_1,\dots,X_n$ are i.i.d.\ Bernoulli$(p)$, $S_n = \sum X_i$, and $\varepsilon>0$, then
\[
P\left(\frac{S_n}{n} \ge p + \varepsilon\right)
\le \exp(-n D(p+\varepsilon \,\|\, p)),
\]
where $D(q\|p)$ is the Kullback--Leibler divergence between Bernoulli$(q)$ and Bernoulli$(p)$.

\subsection{Cram\'er’s Theorem (Large Deviations for Sums)}

Cram\'er's theorem provides a full large deviations principle for normalized sums of i.i.d.\ random variables: for $S_n=\sum_{i=1}^n X_i$, the probabilities
\[
P\left(\frac{S_n}{n} \in A\right)
\]
decay like $\exp(-n \inf_{x\in A} I(x))$, where $I$ is the \emph{rate function} given by the Legendre transform of the log moment generating function:
\[
I(x) = \sup_{t\in\mathbb{R}} \{tx - \Lambda(t)\},
\quad \Lambda(t) = \log E[e^{tX_1}].
\]

This is much sharper than the CLT: it describes the exponential decay of large deviations from the mean.

\begin{example}[Rare-event probabilities for sample mean]
If $X_i$ are daily P\&L shocks with light tails, Cram\'er's theorem can approximate probabilities that the average return over $n$ days exceeds a stressed level. 
The rate function $I(x)$ quantifies how exponentially rare such events are in the i.i.d.\ model, and is central to risk measures like value-at-risk and expected shortfall in large portfolios.

\end{example}

\newpage

%%%%%%%%%%%%%%%%%%%%%%%%%%%%%%%%%%%%%%%%%%%%%%%%%%%%%%%%%%%%
\section{Hints of High-Dimensional Statistics}
\label{sec:high-dim}

\subsection{Curse of Dimensionality and Sparsity}

In high dimensions ($p$ comparable to or larger than $n$), classical asymptotics ($n\to\infty$ with fixed $p$) break down. 
Covariance matrices become singular, and many estimators are not well-defined without additional structure.

A common structural assumption is \textbf{sparsity}: only a small subset of parameters (e.g.\ regression coefficients) are non-zero. 
Regularization methods exploit this to recover signal from limited data.

\subsection{Penalized Estimation (Lasso Idea)}

In linear regression $\bm{y}=X\bm{\beta} + \bm{\varepsilon}$, the Lasso estimator solves
\[
\hat{\bm{\beta}}^{\text{lasso}}
= \arg\min_{\bm{\beta}}
\left\{
\frac{1}{2n}\|\bm{y}-X\bm{\beta}\|_2^2
+ \lambda \|\bm{\beta}\|_1
\right\}.
\]

The $\ell_1$ penalty $\|\bm{\beta}\|_1$ encourages many coefficients to be exactly zero, yielding sparse solutions. 
Concentration inequalities and empirical process theory underpin consistency and oracle inequalities for such estimators in high dimensions.[web:65][web:58]

\subsection{Concentration of Measure}

In high dimensions, Lipschitz functions of many independent variables are often tightly concentrated around their mean or median. 
Results like McDiarmid's inequality and Gaussian concentration show that deviations of order larger than $\sqrt{n}$ are extremely unlikely.

\begin{example}[Norm of a high-dimensional Gaussian vector]
If $\bm{Z}\sim N_d(\bm{0},I_d)$, then $\|\bm{Z}\|_2^2$ is $\chi^2_d$-distributed with mean $d$ and variance $2d$. 
Thus $\|\bm{Z}\|_2$ is concentrated around $\sqrt{d}$. 
For large $d$, the relative fluctuations shrink:
\[
\frac{\|\bm{Z}\|_2}{\sqrt{d}} \xrightarrow{P} 1.
\]
This concentration is fundamental in random matrix theory and high-dimensional geometry, and has strong implications for statistical algorithms in large-$p$ settings.

\end{example}

\newpage

%%%%%%%%%%%%%%%%%%%%%%%%%%%%%%%%%%%%%%%%%%%%%%%%%%%%%%%%%%%%
\section{Where to Go Next}

This advanced course has covered:
\begin{itemize}
    \item Measure-theoretic probability: probability spaces, random variables as measurable maps, Lebesgue integration.
    \item Convergence modes, Borel--Cantelli lemmas, zero--one laws, strong laws, CLTs, and the law of the iterated logarithm.
    \item Conditional expectation as projection, martingales, optional stopping, and convergence theorems.
    \item Discrete-time Markov chains, ergodicity, and continuous-time Brownian motion with a glimpse of stochastic calculus.
    \item Decision-theoretic statistics: risk, Bayes and minimax rules; sufficiency, completeness, UMVU, and information inequalities.
    \item Asymptotic likelihood theory, Bayesian asymptotics, concentration inequalities, large deviations, and a taste of high-dimensional methods.
\end{itemize}

Natural next steps include:
\begin{itemize}
    \item Full courses in stochastic calculus and stochastic differential equations, with applications to finance and control.
    \item Advanced mathematical statistics: empirical processes, semiparametric theory, nonparametric Bayes.
    \item High-dimensional probability and random matrix theory.
    \item Specialized topics: time series (ARMA, GARCH), survival analysis, spatial statistics, causal inference, and modern machine learning theory.
\end{itemize}

Graduate syllabi in advanced probability and mathematical statistics list these directions as core components of a deeper study path.[web:62][web:65][web:66][web:70][web:72]

\end{document}
